\section{Computing as textual re-writes}
\begin{quotation}\textit{
To understand computers fully, and especially to reason about what they
can do and how long it should take, it is important to have a clean model
of what computation is. UNtil that is in place there is a risk that some
of the questions that are raised might have different answers depening on
whether you are thinking in terms in Windows, macOS, Linux or Android or
whether the computer you intend to use is to be programmed in Python, Rust,
C++ or its internal machine language. Obviously there are various and rather
large practical issues that arise because of the above range of choices, but
there remain things that can be said about computing that do not depend
on surface presentation. To get to look at those it is useful to start from
a somewhat impractical and abstract model of computation that starts off
simple enough that everybody can undersand it is very full detail. There are
half a dozen such starting points, and several will appear in later chapters,
but the one followed through on here is based on a couple of textual
transitions on strings of characters. Amazingly it is possible (in not
really many pages) to transition from pretty barbarous raw strings to
use of notations that have recognisable similarity to sensible high
level programming languages.
}\end{quotation}
Back in the 1920s mathematicians were seeking ways of understanding the
very foundations of their disipline. There was work that tried to build
everything on the natural numbers 0,1,\ldots. Others focussed on sets.
A third group started with functions. While this work was all done as
investigations within pure mathematics, these days parts of it can be
viewed computationally. The version here was developed my Moses
Sch\"onfinkel\cite{schonfinkel}.
It starts with just two simple rewrite rules:
\begin{lstlisting}[mathescape=true]
    ((K $x$) $y$) $\Rightarrow$ $x$
    (((S $f$) $g$) $x$) $\Rightarrow$ (($f$ $x$) ($g$ $x$))
\end{lstlisting}
where here the lower case letters are to be viewed as placeholders, so
for instance \verb+((K rabbit) hat)+ will transform into jusr \verb+rabbit+.
This was of starting off is known as the use of ``combinators''.

These two rules set up a computational model if you suppose that an input
``program'' gets presented as just a nicely parenthesized string with only
\verb+S+ and \verb+K+ present (so no rabbits!). Things proceed by having
a search to see if either of the patterns match anywahere in this material,
in which case the indicated transformation is made. The system halts when
there are no more reductions available to perform, and of course it runs
for ever if there are always more cases of the left hand sides of the two
rules however long you go on. There is a technicality that arises because in
lengthy inputs there may be an opportunity to choose between several uses
of the two rules. The traditional and proper resolution to this is that the
leftmost-outermost transition should always be made first. Following that
guidance in fact arranges that if the calculation has any change of
terminating then it will.

This pair of rules starts out by looking odd and not really very powerful.
BUt using a succesion of transformations they can implausibly easily be
shows to relate to sensible computing using reasonable notation.

The adjustment starts with three bits of minor trickery that render everything
a lot easier to work with:
\begin{enumerate}
\item For the sake of the human reader it is very helpful to avoid the need
to write in all the parentheses. This ia achieved sith insistence that
when items are presented one after the other they are to be considered
as grouped to the left. With this understanding the two rules can be written
as
\begin{lstlisting}[mathescape=true]
    K $x$ $y$ $\Rightarrow$ $x$
    S $f$ $g$ $x$ $\Rightarrow$ $f$ $x$ ($g$ $x$)
\end{lstlisting}
and while with as small a case as this the improvement is not overwhelming,
for larger cases the clarity end up hugely better.
\item For writing a program that applies the rules the use of a string
representation and load of brackets would be clumsy. It is way better
to represent the transformations as ones of binary trees:
\begin{lstlisting}[mathescape=true]
           .
          / \
         .   $y$    $\Rightarrow$  $x$
        / \
       K   $x$

           o
          / \
         .   $x$    $\Rightarrow$                       o
        / \                            /  \
       .   $g$                           .    .
      / \                            / \  / \
     S   $f$                           $f$   $x$   $g$   $x$

\end{lstlisting}
where if $x$ is a further tree structure the result from \verb+S+ will
have two references to it, so that the single copy of it is
shared.
It is also reasonable that the node in the tree shown as \verb+o+ gets
overwritten while creating the result-tree from \verb+S+. If that is sone
and two higher level part of the tree share references to it then when the
second one gets there it will find that the required transformation has
already been done!
In this tree interpretation of things the ``leftmost outermost'' location
can generally be found be just scanning down the left spine of the tree
until either and \verb+S+ or \verb+K+ is found.
\item Going back to making life easier for the human reader, it is good to
allow a naming scheme after the style
\begin{verbatim}
    let I = S K K;
\end{verbatim}
that gives names to various items. These names are not thought of as in any
way part of the computation model -- before one starts applying tow two
rules all names can be replaced with what they stand for so that the
computation really only involves the primitive symbols \verb+S+ and
\verb+K+. But the names really make explanations a lot smoother!
\end{enumerate}

Well the above explanation introduced a symbol \verb+I+. So it is
proper to wonder what its properties are, The folliwing shows the
how use of it proceeds, starting with simply replacing the symbol with
its expansion
\begin{lstlisting}[mathescape=true]
    I $a$
    S K K $a$
    K $a$ (K $a$)
    $a$
\end{lstlisting}
and this suggestes that \verb+I+ acts like an identity operator.
In a while it will be convenient to introduce a collection of additional
symbols with e.g.\ names like \verb+B+ and \verb+C+ that can be viewed as
nothing more than shorthand for various mixes of \verb+S+ and \verb+K+ but
that make everything a lot easier to use. But before that it is good to
introduce the most basic parts of the programming language that can be
built on top of this foundation.

The most important part of it is the function definition, and this
is really just an extension of the naming scheme already explained.
For this a function that will take one argument would be written
as for instance
\begin{verbatim}
    let applyToSelf f = f f;
end{verbatim}
At the moment we can not use anything other that the formal argument to the
function and the known names \verb+S+, \verb+K+ \nd \verb+I+ in the body
of the function, so the example here just says ``treat f as a function and
give it itself as an argument''. Aha this refers back to the quine section
which may suggest that although this seems a bit odd there may be
circumstances in which it could arise and be sort of useful. But kust now
it is nothing more than an example.

The challenge is to expand this so that there is just a single symbol to
the left of the equals sign, so it is nothing more frightening than
than the existing scheme for naming things. This can be done unexpctedly
easily via a notation (known as ``abstraction'') that just moves the
format argument over to the right hand side:
\begin{verbatim}
    let applyToSelf = [f] (f f);
end{verbatim}
and then uses three easy rules to get rid of the square-bracketed prefix:
\begin{lstlisting}[mathescape=true]
    [$x$] $x$ $\Rightarrow$ I
    [$x$] $y$ $\Rightarrow$ K $y$
    [$x$] $f$ $g$ \$Rightarrow$ S ([$x$] $f$) ([$x$] $g$)
\end{lstlisting}
where in this presentation $y$ is an literal value not equal to $x$, and in
the rules that introduces \verb+S+ it is understodd that the two little
abstractions (one of $f$ and one on $g$) themselves get expanded out.
These three little rules are in many respects just the cmplements of the
ones that defined \verb+S+, \verb+K+ and \verb+I+ so how can anybody
describe them as ``difficult''.

So what about functions with two or more arguments?
\begin{lstlisting}[mathescape=true]
   let f $a$ $b$ $c$ = ...    $\Rightarrow$
   let f $a$ $b$ = [$c$] ...
\end{lstlisting}
and the arguments move over onto the right hand side one at a time. The
consequence of all this is that it now proper to start writing little
sequences of function definitions in a fairly human readable form
safe in the knowledge that their behaviour is fully explanable in terms of
out two simple initial rewrite rules. But at present we have no
data types or useful control structures. Well that is going to be east to
fix:
\begin{verbatim}
   let true a b = a;
   let false a b = b;
   let if p x y = p x y;
\end{verbatim}
intoduced the truth values and in ``if'' operator. It is easy to verify that
for instance \verb+if true A B+ yields \verb+A+. In many programming
contexts one would worry that in the process \verb+B+ would have been
inspacted and it it was an expression that could abort (or get stuck in an
unending computation) there would be trouble. It turns out that the
outermost-leftmost'' rule has the effect of arranging that in this case and
others like it the unwabted value \verb+B+ is not inspected and ithus iti
can not cause trouble. So for now on we can write code that uses \verb+if+
although there are not really any tests available to use it with!

The next big trick here will be to incorporate a modelling of numbers.
This is doneviewing a number $N$ as basically an instruction to
perform some operation $N$ times. So we will have
\begin{lstlisting}[mathescape=true]
   let zero $f$ $x$ = $x$;
   let one $f$ $x$ = $f$ $x$;
   let two $f$ $x$ = $f$ ($f$ $x$);
   let three $f$ $x$ = $f$ ($f$ ($f$ $x$));
\end{lstlisting}
and so on. Basically $N f$ denotes the function $f$ composed with itself $N$
times. Two fundamental things we need are a way of testing of such
a number is zero and a way of adding one:
\begin{lstlisting}[mathescape=true]
   iszero $N$ = $N$ (K false) true;
   add1 $N$ $f$ $x$ = $f$ ($N$ $f$ $x$);
\end{lstlisting}
where each of those definitions are so short and simple that one can just
;check them by seeing what happens when theya are given representative
arguments. In manners that may or may not be immediately obvious ;it is
then easy to write
\begin{lstlisting}[mathescape=true]
   let plus $M$ $N$ $f$ $a$ = $M$ $F$ ($N$ $f$ $a$)
   let times $M$ $N$ $f$ = $M$ ($N$ $f$)
   let power $M$ $N$ = $N$ $M$;
\end{lstlisting}
and even though it may take a few moments and a bit of experimentation to
become convinced that the code for \verb+power+ will raise one number ti
a power as specified by the other bobody can dispute how compact it is.

Those concerned with practicality can now take the view that when they write
programs that involve these names they might as well use ordinary notation
such as 0, 1, 2, 3 for the numbers, \verb/+/ and \verb/*/ for addition
and multiplication (using these as infix operators), and express their
conditionals using syntax with words \verb+if+, \verb+then+ and \verb+else+.
When it comes to pedantic reasoning to prove properties of the code they
write it would be easy to unwind all the way back to express things in terms
of the ultimate basics.

AT this stage it is perhaps useful to step back and think how much has
been achieved in a really modest number of lines of definitions given
that raw \verb+S+ and \verb+K+ seem to extrememly unpromising. And
also to recognize how much each of those definitions looks like a very
simple fragment of some real program.

The above shows one way to incorporate aritmetic into this world, but
does not (thus far) illustrate subtraction and division. This is
possible in not many more lines of definitions, but it will perhaps be more
interesting here to take a different track.
One of the most basic and hence
useful buidling blocks for data structures is the {\em pair}. This takes
two values and glues them together such that at a later stage it will be
possible to extract each component. If you have pairs it will be possible
to chain them as lists and view the length of the list as an alternative
denotation of a number. If one can concatenate lists that corresponds
to addition, and taking the tail of a list subtracts one. So in fact by
looking at data representation we find that a fresh way of modelling arithmetic
emerges automatically -- one where subtracting 1 is easy to achieve by
where multiplication is more of a stress. Obviously as well as
using pairs to model numners they can be used to build up vectors and
trees and all manner of other data-structures. The caveat for now is that
the structures are read-only so if you though you wanted to overwrite
some part of one you will in fact need to make an edited copy that
reflects the desired change. Well this sounds powerful, so is it hard to
achieve? Well here is a function \verb+pair+ that create a pair
and functions \verb+left+ and \verb+right+ to extract the two
components:
\begin{lstlisting}[mathescape=true]
   let (pair $a$ $b$) $f$ = $f$ $a$ $b$;
   let left $p$ = $p$ true;
   let right $p$ = $p$ false;
\end{lstlisting}
In this presentation there are some extra parentheses in the
exposition of \verb+pair+ just to stress that it will usually be
called with only two arguments. In this world calling a function with
fewer arguments than had been shown in the textual definition is not
to be seen as a problem. All that happens is that one gets left with an object that
can sit around only partially evaluated, but that is awaiting the remaining argument
or arguments to complete its work.

The above version of \verb+pair+ is good as far as it goes, but if
it is to be used to build lists or other structures whose precise shape
is not known in advance it may be necessary to tag each little buillding
block with a boolean value marking whether it is a terminator or not:
\begin{lstlisting}[mathescape=true]
   let emptylist = pair true anything;
   let nonemptylist $head$ $tail$ = pair false (pair $head$ $tail$);
   let isempty $l$ = left $l$
   let first $l$ = left (right $l$);
   let rest $l$ = right (right $l$);
\end{lstlisting}
and again these code fragments are concise, not that obviously
terribly inefficient and they still look as if they are written
for the benefit of humans not just machines. Just as once numbers had been
introduced one could expand the shorthand notation to use 0, 1, 2 and 3 here
one might well permit some notation such as \verb+(A,B)+ for a raw pair
and \verb+{A,B,C,D}+ for a list. The level of syntactic sugar makes working
in this combinator-based world almost attractive!

One final but rather large trick remains. So far the \verb+let+
definitions can not include forward references and functions can not
be defined in terms of themselves (i.e.\ recursively). If that limitation
can be fixed we will really be on the way to writing any program we want
and being able to reason about it with extreme precision by relating
it back to the foundations that \verb+S+ and \verb+K+ provide. The
scheme that achieved this involved inventing a symbol that is traditionally
called $Y$ such that it behaves as if it had been defined using
\begin{lstlisting}[mathescape=true]
   let Y $f$ = $f$ (Y $f$); 
\end{lstlisting}
where of course that precise formulation not acceptable because it
tries to define \verb+Y+ in terms of itself.
However almost be magic the following will work except that it now
uses a qualifier ``where'' to introduce a local definition.
\begin{lstlisting}[mathescape=true]
   let Y $f$ = G G
      where G $x$ = $f$ ($x$ $x$);
\end{lstlisting}
but still if you follow that definition through it satisfies the behaviour
being sought. So here is one further attempt:
\begin{lstlisting}[mathescape=true]
   let H $f$ $x$ = $f$ ($x$ $x$);
   let Y $f$ = (H $f$) (H $f$);
\end{lstlisting}
and now \verb+Y+ is defined in ways that do not rely on any form of
self or forward reference and so what we have would expand into a
mess of uses of just \verb+S+ and \verb+K+ in a perfectly straightforward
way. Note that the input form ``\verb+Y+ $f$`` (for any $f$ you like to use) does
not just generate itself as its output - it generates something that contains
itself. So if forms a sort of super-quine.

We can now exploit $Y$ by arranging that going forwards any \verb+let+
declaration is expanded a little more, with
\begin{lstlisting}[mathescape=true]
   let something $x$ $y$ = something;
\end{lstlisting}
getting rewritten as
\begin{lstlisting}[mathescape=true]
   let something = Y ([something] ([$x$] ([$y$] something)));
\end{lstlisting}
where the square brackets denote abstraction as used earlier where the
treatment of \verb+let+ was introduced. The consequence of this slightly
more elaborate conversion is that the function \verb+something+ can now
be a recursive one.

Because this book is all about difficulties (and also because this is
not actuallly terribly messy once one has got into the swing of things!)
this transformation is just quoted here, and the manner in which
it lets the function \verb+something+ recurse is left for the reader
to trace through and become convinced by.

There is a great deal more that could be discussed here showing how
to handle pairs of mutually recursive functions, how to get behaviours
that are rather like having assigment statements and arrays that
can have their elements updates, modelling of the \verb+catch+ and
\verb+throw+ exception handling facilities of full scale languages and
almost anything else from the programming world that one can imagine.
IN most cases the steps needed to accomodate the extra capabilities are
not much more messy than those used so far -- the problem is more that
there are quite a few of them and in reality the towers of them tend to
lead to the \verb+S+ and \verb+K+ level rendition getting very bulky.
So there is also a great deal more that could be discussed about
enhancing the undepinning model with new operations of which the first
would be just
\begin{lstlisting}[mathescape=true]
   let B $f$ $g$ $x$ = $f$ ($g$ $x$;
   let C $f$ $x$ $y$ = $f$ $y$ $x$;
\end{lstlisting}
and where the rules for abstraction would spot cases where these could
profitably be used. At the end of that path and with a modest handful of
extra primitives it even turns out that the conversion from the
human-readable source forms of functions to primotive form hardly
expands it at all. As an (unexplained) illustration of this
here is a function for computing a factorial and its converted form:
\begin{verbatim}
fun fact n =
   if n = 1 then 1
   else n * fact(n-1);

Y (B (S (C' if (= 1) 1)) (B (S *) (C B (C - 1))))
\end{verbatim}

What should have emerged from this section is a recognition that
just because something looks as simple and harmless as the
pair of rules
\begin{lstlisting}[mathescape=true]
    ((K $x$) $y$) $\Rightarrow$ $x$
    (((S $f$) $g$) $x$) $\Rightarrow$ (($f$ $x$) ($g$ $x$))
\end{lstlisting}
seems to be there is no certainty that the computational behaviour
described is in any way simple, and in this case pretty well arbitrary
programs with all the challenges of understanding them can be
represnted based on the utterly primitive foundation captures in these
two almost trivial rules.
.
